\section{Introduction.}

Most of the concrete results of this paper were obtained by methods of
lattice field theory. These results make up a first component of a
larger, rather ambitious and quite speculative
program which is defined in the continuum, independently of any
regularization. The outline of the program makes up the bulk of
this section. The intention is to put the presentation of the lattice
results that follows in later sections in perspective. Throughout
the paper, we shall get back to continuum language to emphasize
what structures are, in our view, continuum properties.


Pure four dimensional $SU(N)$ gauge theories are strongly
interacting at large distances and weakly interacting at short
scales. In $SU(3)$ the cross-over between these two regimes is
relatively narrow and there are no good analytic approximation
methods that apply to it; nor do we have a continuum way to
calculate at strong coupling. The cross-over narrows as the number
of colors, $N$, increases.

Here, in a direct continuation of the line of thought of some old
work \cite{old1, gw12}, we report on an improved line of attack on this
problem and take several steps in its direction. The old idea was
that instantons indicate that in the limit of infinite $N$, at a
particular sharply defined scale, there is a phase transition, and
weak and strong coupling physics can be connected at that point.
This idea has not been successfully implemented because too many
parts were either wrong or missing. Much has been learned in the
interim period and we would like to try again. Let us now briefly
overview the new ingredients of our approach:

While the guess that the cross-over regime
collapses to a point and becomes a phase transition at
infinite $N$ is old, we
understand now that the scale at which the transition occurs is
not a universal quantity that applies to all observables
simultaneously; nor does it necessarily affect any local
observables: it is not seen in the free energy density for
example.
The most interesting large $N$ transitions do not have
to be ``bulk'' phase transitions.

Recent lattice work on large $N$ has produced such ``bulk''
transitions at global
critical scales in the planar limit of QCD on a Euclidean torus,
but the associated transitions take the system out of the zero
temperature, infinite volume regime and depend on the global
topology of space-time. This is not going to directly explain how
a few particle high energy process hadronizes, for example.

Recent lattice work has also indicated that
specific non-local observables undergo large
$N$ phase transitions as they are dilated \cite{bulk2,knn} and this
happens while the Euclidean space-time volume stays infinite in
all directions. The critical scales are observable dependent; one
can imagine associating with these observables an entire family of
dilated images and somewhere along the dilation axis the large $N$
transition takes place.

If there were a string equivalent to
planar QCD, one would guess that the dilation axis is somehow
related to a fifth coordinate as has been often hypothesized
\cite{poly3,malda4}. A meaningful higher dimensional formulation
presupposes a local definition of the theory in the extended
space; we shall have something to add to this speculation at the
end of the paper, but it will be vague.

Another new ingredient we add to the old program is the modern
understanding of universal features of certain large $N$
transitions in matrix models. Although, on the face of it, there
is no concept of locality in the index space associated with the
symmetry group, the models admit a
't Hooft double line expansion for two-index representations
classified by the genus of two
dimensional surfaces. A concept of two-dimensional locality
applies to the physics on these surfaces and the associated
Liouville type models represent universality classes of finite,
but varying degrees of stability associated with the large $N$
transitions \cite{doublesca5}. Most of these models are exactly
solvable.

We hope that the transitions we shall exhibit belong to some such
universality class. We have a guess for the relevant universality
class and it has an exactly solvable matrix model representative.
This means that at finite but large enough
$N$, when the transition dissolves into a cross-over, there is a
description of well chosen observables undergoing the cross-over
in terms of universal scaling functions with a dependence on a
limited number of non-universal parameters that typically come in
as boundary conditions on the equations determining the universal,
known scaling functions.

It is through these parameters that the
cross-over regime connects to the weak and strong coupling regimes
on either side. The four dimensional dilation parameter, in the
vicinity of its critical value plays the role of a
coupling that turns critical in the relevant version of matrix model
description of the large $N$ transition. By the mechanism of a
(double) scaling limit the dilation parameter
gets connected to a scale in
the two dimensional physics that codifies the universal features
of the large $N$ transition. This connection will involve a
nontrivial power and an associated dimensionful parameter.

The weak and strong coupling regimes are governed by their own
universal behavior, field theoretical on the weak side, and (we
guess) effective string like on the strong coupling side. Thus,
the free, nonuniversal, parameters needed by
each universal description, field
theoretic, large $N$ and string-like are matched to each other as
the operator is dilated, and several universality mechanisms
coexist in the same system, controlling separate scale ranges.

If the above scenario holds at $N=\infty$, we can apply it to pure
$SU(3)$ gauge theory, and, hopefully provide a
calculation of the string tension in terms of $\Lambda_{\rm QCD}$
which is accurate to order $\frac{1}{ N^2}\approx 10$ percent. For
a physicist, this might be of more value even  than a mathematical
proof of confinement. We are far from realizing this goal.

The first hurdle to overcome is to find good observables; we claim
that the right question is not whether or not there is a large $N$
phase transition in general, but rather, can one find good
observables which undergo such a transition and can be controlled
(by different theories) at very short and at very long distance
scales respectively. Local observables are unlikely to undergo the
transition. Therefore the cherished analyticity properties in
four dimensional Euclidean momentum space stay protected,
maintaining the basic unitarity and covariance properties of the
field theoretical description of the physics.

In this paper we define a
class of smeared Wilson loop observables on the lattice with the
intent that they have finite continuum limits and obey the
following properties:
\begin{itemize}
\item It is correct to think about the parallel
transport round the loop as being realized by an $SU(N)$ matrix.
This statement has to hold after the renormalization of the theory
and of the operator have been carried out.
\item When we scale the loop, a phase
transition takes place in its eigenvalue distribution if
$N=\infty$.
\item The smeared loop operator is calculable in
perturbation theory when it is scaled down.
\item Its shape dependence
is well described by an effective string theory when it is scaled
up.
\item It is possible to see by continuum non-perturbative
methods (e. g. instantons) that the phase transition is approached
when the loop is dilated up from being very small, so one can hope
to calculate the matching at the weak coupling end of the
cross-over.
\end{itemize}
At the large scale end one also needs a way to
calculate a matching, and this is presumed to work by the large
scale end of the cross-over having some string description (coming
from the relevant universal matrix model) which can be joined
smoothly to the right effective four dimensional effective
string theory. That effective string theory
could be, for example, similar
to \cite{ps6}. This we have not yet established, and is left for
future work. Nevertheless, on the basis of the results
we do have, we believe to have found
non-local observables which are ``good'' enough.

In its essence the calculation we have
outlined is an attempt to extend the
logic of effective field theory \cite{effft7} calculations to the
strong-weak cross-over in pure gauge QCD. Also, the program is in
the same spirit as the old CDG approach to instantons
\cite{cdg75} where they tried to use instanton
gas methods to set the parameters in an effective bag-model
description of hadrons.

If our scenario works for pure glue, it
would be possible to re-introduce fermions at a later stage, at
least as probes with no back reaction and see whether one can
construct hadronic observables containing fermionic fields that
also undergo large $N$ transitions. In particular it would be
interesting to find a
similar calculational program
for spontaneous chiral symmetry breaking.
One could check this numerically using exactly chiral
fermions~\cite{overlap}.

In the next section we introduce our observables on the lattice
and present our numerical results in support of the claim that our
observables have a continuum limit and undergo large $N$
transitions under dilation. In subsequent sections we discuss a
continuum view of the lattice construction and our guess for the
large $N$ universality class. The presentation
becomes by necessity progressively more speculative. Ultimately,
we hope to reduce the role of lattice simulations to a strictly
qualitative one, providing support for the assumptions about the
various universality classes that control the different regimes
our calculation needs to connect.
The calculation in itself would be entirely
analytical. It goes without saying that any analytical step would
need to be checked against numerical work, as this is unchartered
territory.
